\documentclass[11pt]{article}
\usepackage[margin=1in]{geometry}
\usepackage[T1]{fontenc}
\usepackage[utf8]{inputenc}
\usepackage{lmodern}
\usepackage{microtype}
\usepackage{amsmath,amssymb,amsthm,mathtools}
\usepackage{bm}
\usepackage{enumitem}
\usepackage{hyperref}
\hypersetup{colorlinks=true,linkcolor=blue,urlcolor=blue,citecolor=blue}
\allowdisplaybreaks

\newtheorem{theorem}{Theorem}
\newtheorem{proposition}[theorem]{Proposition}
\newtheorem{corollary}[theorem]{Corollary}
\newtheorem{lemma}[theorem]{Lemma}
\newtheorem{remark}[theorem]{Remark}
\newtheorem{assumption}{Assumption}
\newtheorem{assumptionprime}{Assumption}
\renewcommand{\theassumptionprime}{\theassumption$'$}
\theoremstyle{definition}
\newtheorem{definition}{Definition}
\theoremstyle{plain}

\newcommand{\E}{\mathbb{E}}
\newcommand{\1}{\mathbf{1}}
\newcommand{\plim}{\mathop{\mathrm{plim}}}
\newcommand{\R}{\mathbb{R}}

\title{Design-based interpretation of Poisson regression}
\author{Georgy Kalashnov, Lihua Lei}
\date{\today}

\begin{document}
\maketitle

\section{Setup}
\label{sec:setup}

Let $(Y_i,W_i,X_i)_{i=1}^n$ be an i.i.d. sample, where $W_i\in\{0,1\}$ and $Y_i=Y_i(W_i)\ge 0$. We suppress the unit subscript below and will denote a generic observation by $(Y,W,X)$.

\begin{assumption}[Integrability]
\label{as:integrability}
$\E[Y(0)+Y(1)]<\infty$.
\end{assumption}

Let $e(X):=\Pr(W=1\mid X)$.

The estimand $\tau^\star$ is the coefficient on $W$ in the population Poisson regression with log link, which fits the conditional mean
\begin{equation}
\exp\{h(X)+\tau W\}=\mu_0(X)e^{\tau W},
\qquad
\mu_0(X):=\exp\{h(X)\}\ge 0 .
\label{eq:loglink}
\end{equation}

Let $\mathbb P$ be the law of $(Y(0),Y(1),W,X)$.  For $\tau\in\R$ let $\E_\tau$ denote expectation under the tilted measure $\mathbb P_\tau$ with
\begin{equation}
\frac{d\mathbb P_\tau}{d\mathbb P}:=\frac{e^{\tau W}}{\E[e^{\tau W}]},
\label{eq:tilt}
\end{equation}
expectations without a subscript being under $\mathbb P$.  Let $C_\tau(X):=\E[e^{\tau W}\mid X]=(1-e(X))+e(X)e^{\tau}$; since $W\in\{0,1\}$, $\E[e^{\tau W}]=\E[C_\tau(X)]$ lies in $[\min(1,e^{\tau}),\max(1,e^{\tau})]$.  Under $\mathbb P_\tau$ the conditional law of $W$ given $X$ is
\begin{equation}
\Pr\nolimits_\tau(W=w\mid X)=\frac{e^{\tau w}\Pr(W=w\mid X)}{C_\tau(X)},
\qquad w\in\{0,1\}.
\label{eq:condtilt}
\end{equation}
The density in \eqref{eq:tilt} depends on $W$ only, so $\mathbb P_\tau$ leaves the conditional law of $(Y(0),Y(1))$ given $(W,X)$ unchanged.

For $\tau\in\R$ let
\begin{equation}
\pi_\tau(X):=\E_\tau[W\mid X]=\frac{e(X)e^{\tau}}{C_\tau(X)}
\label{eq:defpi}
\end{equation}
be the \emph{tilted propensity score}, the second equality by \eqref{eq:condtilt}.

The purpose of this paper is to find conditions, under which we can interpret the estimand of the Poisson regression as an average of potential outcomes, weighted by the ex-post or the ex-ante weight.

\begin{definition}[Ex-post weight]
\label{def:expost}
\begin{equation}
\psi_\tau(W,X):=\bigl(W-\pi_\tau(X)\bigr)\1\{W>0\}=\bigl\{1-\pi_\tau(X)\bigr\}W .
\label{eq:defexpost}
\end{equation}
\end{definition}

\begin{definition}[Ex-ante weight]
\label{def:exante}
\begin{equation}
\phi_\tau(X):=\E\bigl[\psi_\tau(W,X)\mid X\bigr]=\frac{e(X)(1-e(X))}{C_\tau(X)} .
\label{eq:defexante}
\end{equation}
\end{definition}

By such an interpretation we mean that $\tau^\star$ solving the population Poisson first-order conditions \eqref{eq:condFOC}--\eqref{eq:treatFOC} satisfies
\[
\tau^\star=\log\frac{\E\bigl[\omega\,Y(1)\bigr]}{\E\bigl[\omega\,Y(0)\bigr]},
\qquad
\omega\in\bigl\{\psi_{\tau^\star}(W,X),\ \phi_{\tau^\star}(X)\bigr\}.
\]

\section{Design-Based Assumption}

\begin{assumption}[Design-based]
\label{as:design}
$\E\bigl[W\bigm| Y(0),Y(1),X\bigr]=\E\bigl[W\bigm| X\bigr]=e(X)$ almost surely.
\end{assumption}

Under Assumption \ref{as:design} the conditional law of $W$ given $(Y(0),Y(1),X)$ equals its conditional law given $X$, so for every $\tau\in\R$
\begin{equation}
\E\bigl[e^{\tau W}\bigm| Y(0),Y(1),X\bigr]=C_\tau(X),
\qquad
\E_\tau\bigl[W\bigm| Y(0),Y(1),X\bigr]=\E_\tau\bigl[W\bigm| X\bigr]=\pi_\tau(X).
\label{eq:designtilt}
\end{equation}

\begin{proposition}
\label{prop:reduction}
Suppose Assumptions \ref{as:integrability} and \ref{as:design} hold. Let $\tau^\star\in\R$ and a measurable function $\mu_0^\star(X)\ge 0$ satisfy the population first-order conditions
\begin{align}
\E\big[Y-\mu_0^\star(X)e^{\tau^\star W}\mid X\big] &= 0 \quad \text{a.s.},
\label{eq:condFOC}
\\
\E\big[W\{Y-\mu_0^\star(X)e^{\tau^\star W}\}\big] &= 0.
\label{eq:treatFOC}
\end{align}
Then
\begin{equation}
\E\Big[\phi_{\tau^\star}(X)\bigl\{Y(1)-e^{\tau^\star}Y(0)\bigr\}\Big]=0.
\label{eq:pseudotrue}
\end{equation}
If in addition $\E[\phi_{\tau^\star}(X)Y(0)]>0$, then
\begin{equation}
\tau^\star=\log\frac{\E[\phi_{\tau^\star}(X)\,Y(1)]}{\E[\phi_{\tau^\star}(X)\,Y(0)]} .
\label{eq:expostratio}
\end{equation}
\end{proposition}

\begin{proof}
We first change measure from $\mathbb P$ to $\mathbb P_{\tau^\star}$.  Since $Y-\mu_0^\star(X)e^{\tau^\star W}=e^{\tau^\star W}\bigl\{e^{-\tau^\star W}Y-\mu_0^\star(X)\bigr\}$, $\E[e^{\tau^\star W}]>0$ and $C_{\tau^\star}(X)>0$, \eqref{eq:tilt} and \eqref{eq:condtilt} turn \eqref{eq:condFOC}--\eqref{eq:treatFOC} into
\begin{align}
\E_{\tau^\star}\big[e^{-\tau^\star W}Y-\mu_0^\star(X)\bigm| X\big] &= 0 \quad \text{a.s.},
\label{eq:tiltedcondFOC}
\\
\E_{\tau^\star}\big[W\{e^{-\tau^\star W}Y-\mu_0^\star(X)\}\big] &= 0.
\label{eq:tiltedtreatFOC}
\end{align}
Multiplying \eqref{eq:tiltedcondFOC} by $\pi_{\tau^\star}(X)$, using the law of iterated expectations and subtracting from \eqref{eq:tiltedtreatFOC},
\begin{equation}
\E_{\tau^\star}\Bigl[(W-\pi_{\tau^\star}(X))\bigl\{e^{-\tau^\star W}Y-\mu_0^\star(X)\bigr\}\Bigr]=0 .
\label{eq:firstsubtraction}
\end{equation}
Since $\mu_0^\star(X)$ is a function of $X$ and $\E_{\tau^\star}[W-\pi_{\tau^\star}(X)\mid X]=0$, \eqref{eq:firstsubtraction} gives
\begin{equation}
\E_{\tau^\star}\bigl[(W-\pi_{\tau^\star}(X))\,e^{-\tau^\star W}Y\bigr]
=
\E_{\tau^\star}\bigl[(W-\pi_{\tau^\star}(X))\,\mu_0^\star(X)\bigr]
=0 .
\label{eq:residualized}
\end{equation}

Since $W\in\{0,1\}$,
\[
\bigl(W-\pi_{\tau^\star}(X)\bigr)e^{-\tau^\star W}Y
=
e^{-\tau^\star}\bigl\{1-\pi_{\tau^\star}(X)\bigr\}W\,Y(1)-\pi_{\tau^\star}(X)(1-W)\,Y(0).
\]
By \eqref{eq:designtilt},
\[
\E_{\tau^\star}\bigl[(W-\pi_{\tau^\star}(X))\,e^{-\tau^\star W}Y\bigm| Y(0),Y(1),X\bigr]
=
\pi_{\tau^\star}(X)\bigl\{1-\pi_{\tau^\star}(X)\bigr\}\bigl\{e^{-\tau^\star}Y(1)-Y(0)\bigr\}.
\]

We now change measure back to $\mathbb P$:
\begin{align*}
&\E[e^{\tau^\star W}]\,\E_{\tau^\star}\Bigl[\pi_{\tau^\star}(X)\bigl\{1-\pi_{\tau^\star}(X)\bigr\}\bigl\{e^{-\tau^\star}Y(1)-Y(0)\bigr\}\Bigr]
\\
&\qquad=
\E\Bigl[e^{\tau^\star W}\,\pi_{\tau^\star}(X)\bigl\{1-\pi_{\tau^\star}(X)\bigr\}\bigl\{e^{-\tau^\star}Y(1)-Y(0)\bigr\}\Bigr]
\\
&\qquad=
\E\Bigl[C_{\tau^\star}(X)\,\pi_{\tau^\star}(X)\bigl\{1-\pi_{\tau^\star}(X)\bigr\}e^{-\tau^\star}\bigl\{Y(1)-e^{\tau^\star}Y(0)\bigr\}\Bigr]
\\
&\qquad=
\E\Bigl[\phi_{\tau^\star}(X)\bigl\{Y(1)-e^{\tau^\star}Y(0)\bigr\}\Bigr] ,
\end{align*}
the second equality by \eqref{eq:designtilt} and the third since $C_\tau(X)\pi_\tau(X)\{1-\pi_\tau(X)\}=e^{\tau}\phi_\tau(X)$.  By \eqref{eq:residualized} and the law of iterated expectations the left-hand side equals 0, which gives \eqref{eq:pseudotrue}, hence
\[
\E\bigl[\phi_{\tau^\star}(X)Y(1)\bigr]=e^{\tau^\star}\,\E\bigl[\phi_{\tau^\star}(X)Y(0)\bigr],
\]
which gives \eqref{eq:expostratio} since $\E[\phi_{\tau^\star}(X)Y(0)]>0$.
\end{proof}

The same conclusions hold with a restricted control function, under a design-based condition on the propensity score alone.

Let $R(X)\in\R^K$ be the vector of included controls, containing a constant.

\begin{assumptionprime}[Design-based, restricted controls]
\label{as:designspan}
\begin{equation}
\pi_{\tau^\star}(X)\in \mathrm{span}(R(X))
\qquad \text{almost surely.}
\label{eq:spancondition}
\end{equation}
\end{assumptionprime}

\begin{corollary}
\label{cor:finite}
Suppose Assumptions \ref{as:integrability}, \ref{as:design} and \ref{as:designspan} hold and let $(\beta^\star,\tau^\star)$ solve the population Poisson moment conditions
\begin{equation}
\E\Big[R(X)\bigl\{Y-\mu_0^\star(X)e^{\tau^\star W}\bigr\}\Big]=0,
\qquad
\E\Big[W\bigl\{Y-\mu_0^\star(X)e^{\tau^\star W}\bigr\}\Big]=0,
\label{eq:finiteFOC}
\end{equation}
where $\mu_0^\star(X):=\exp\{R(X)'\beta^\star\}$.  Then the conclusions of Proposition \ref{prop:reduction} hold.  If the sample Poisson estimator $(\hat\beta,\hat\tau)$ is consistent, its limit $\tau^\star$ is characterized by \eqref{eq:pseudotrue} and \eqref{eq:expostratio}.
\end{corollary}

\begin{proof}
We first change measure from $\mathbb P$ to $\mathbb P_{\tau^\star}$: as in the proof of Proposition \ref{prop:reduction}, \eqref{eq:tilt} turns \eqref{eq:finiteFOC} into
\[
\E_{\tau^\star}\Bigl[R(X)\bigl\{e^{-\tau^\star W}Y-\mu_0^\star(X)\bigr\}\Bigr]=0,
\qquad
\E_{\tau^\star}\Bigl[W\bigl\{e^{-\tau^\star W}Y-\mu_0^\star(X)\bigr\}\Bigr]=0 .
\]
By Assumption \ref{as:designspan} there is $\lambda^\star\in\R^K$ with $\pi_{\tau^\star}(X)=R(X)'\lambda^\star$ almost surely, so the first condition gives
\[
\E_{\tau^\star}\Bigl[\pi_{\tau^\star}(X)\bigl\{e^{-\tau^\star W}Y-\mu_0^\star(X)\bigr\}\Bigr]
=\lambda^{\star\prime}\,\E_{\tau^\star}\Bigl[R(X)\bigl\{e^{-\tau^\star W}Y-\mu_0^\star(X)\bigr\}\Bigr]=0 .
\]
Subtracting this from the second condition gives \eqref{eq:firstsubtraction}, and the proof of Proposition \ref{prop:reduction} applies from this point unchanged.
\end{proof}

\begin{remark}
Ex-ante weights are expected values of ex-post weights:
\[
\E\big[\{1-\pi_{\tau^\star}(X)\}W \,\big|\, X\big] = \phi_{\tau^\star}(X) .
\]

Additionally, this is also true under linearity constraints: in the restricted-control-function setting below, if $a^\star(X)=\pi_{\tau^\star}(X)$, which is valid whenever $\pi_{\tau^\star}(X)\in\mathrm{span}(R(X))$ as in Corollary \ref{cor:finite}, then $\E\big[\{1-a^\star(X)\}W \,\big|\, X\big] = \phi_{\tau^\star}(X)$ as well.
\end{remark}

\section{Model-Based Assumption}
Let $R(X)\in\R^K$ be the vector of included controls, containing a constant, and let $(\beta^\star,\tau^\star)$ solve the population Poisson moment conditions
\begin{equation}
\E\Big[R(X)\bigl\{Y-\mu_0^\star(X)e^{\tau^\star W}\bigr\}\Big]=0,
\qquad
\E\Big[W\bigl\{Y-\mu_0^\star(X)e^{\tau^\star W}\bigr\}\Big]=0,
\label{eq:poissonFOC}
\end{equation}
where $\mu_0^\star(X):=\exp\{R(X)'\beta^\star\}$.

\begin{assumption}[Correctly specified control function]
\label{as:spanmeanindep}
$\E[Y(0)\mid W,X]=\mu_0^\star(X)$.
\end{assumption}

\begin{proposition}
\label{prop:expostfinite}
Suppose Assumption \ref{as:integrability} and Assumption \ref{as:spanmeanindep}\footnote{The proof uses Assumption \ref{as:spanmeanindep} only through $\E_{\tau^\star}\bigl[\bigl(W-a^\star(X)\bigr)\bigl\{Y(0)-\mu_0^\star(X)\bigr\}\bigr]=0$, which is therefore sufficient.} hold, $\E_{\tau^\star}[R(X)R(X)']$ is nonsingular, and $a^\star\in\mathrm{span}(R(X))$ satisfies
\begin{equation}
\E_{\tau^\star}\Big[R(X)\bigl\{W-a^\star(X)\bigr\}\Big]=0 .
\label{eq:tiltedprojection}
\end{equation}
Then, with $\psi^\star(W,X):=\bigl\{1-a^\star(X)\bigr\}W$,
\begin{equation}
\E\Big[\psi^\star(W,X)\bigl\{Y(1)-e^{\tau^\star}Y(0)\bigr\}\Big]=0 .
\label{eq:expostfinitemoment}
\end{equation}\footnote{Under this stronger Assumption \ref{as:spanmeanindep} the proof does not use \eqref{eq:tiltedprojection}, so \eqref{eq:expostfinitemoment} holds with $a^\star$ replaced by any $a\in\mathrm{span}(R(X))$.  The same is true of the least squares specification: if $(\beta^\star,\gamma^\star)$ solve $\E[R(X)\{Y-\beta^\star W-R(X)'\gamma^\star\}]=0$ and $\E[W\{Y-\beta^\star W-R(X)'\gamma^\star\}]=0$, and $\E[Y(0)\mid W,X]=R(X)'\gamma^\star$, then $\E[\{1-a(X)\}W\{Y(1)-Y(0)-\beta^\star\}]=0$ for every $a\in\mathrm{span}(R(X))$.}
If in addition $\E[\psi^\star(W,X)Y(0)]>0$, then
\[
\tau^\star=\log\frac{\E[\psi^\star(W,X)\,Y(1)]}{\E[\psi^\star(W,X)\,Y(0)]} .
\]
\end{proposition}

\begin{proof}
We first change measure from $\mathbb P$ to $\mathbb P_{\tau^\star}$.  Since $Y-\mu_0^\star(X)e^{\tau^\star W}=e^{\tau^\star W}\bigl\{e^{-\tau^\star W}Y-\mu_0^\star(X)\bigr\}$ and $\E[e^{\tau^\star W}]>0$, \eqref{eq:tilt} turns \eqref{eq:poissonFOC} into
\begin{equation}
\E_{\tau^\star}\Bigl[R(X)\bigl\{e^{-\tau^\star W}Y-\mu_0^\star(X)\bigr\}\Bigr]=0,
\qquad
\E_{\tau^\star}\Bigl[W\bigl\{e^{-\tau^\star W}Y-\mu_0^\star(X)\bigr\}\Bigr]=0 .
\label{eq:tiltedFOC}
\end{equation}
Since $a^\star=R(X)'\lambda^\star$ for some $\lambda^\star\in\R^K$, the first condition in \eqref{eq:tiltedFOC} gives
\[
\E_{\tau^\star}\Bigl[a^\star(X)\bigl\{e^{-\tau^\star W}Y-\mu_0^\star(X)\bigr\}\Bigr]
=\lambda^{\star\prime}\,\E_{\tau^\star}\Bigl[R(X)\bigl\{e^{-\tau^\star W}Y-\mu_0^\star(X)\bigr\}\Bigr]=0 .
\]
Subtracting this from the second condition in \eqref{eq:tiltedFOC},
\begin{equation}
\E_{\tau^\star}\Bigl[(W-a^\star(X))\bigl\{e^{-\tau^\star W}Y-\mu_0^\star(X)\bigr\}\Bigr]=0 .
\label{eq:expostsubtraction}
\end{equation}

Since $W\in\{0,1\}$,
\begin{equation}
e^{-\tau^\star W}Y-\mu_0^\star(X)
=
W\bigl\{e^{-\tau^\star}Y(1)-Y(0)\bigr\}+\bigl\{Y(0)-\mu_0^\star(X)\bigr\},
\label{eq:expostidentityfinite}
\end{equation}
so multiplying \eqref{eq:expostidentityfinite} by $W-a^\star(X)$ and using $(W-a^\star(X))W=\psi^\star(W,X)$, \eqref{eq:expostsubtraction} becomes
\begin{equation}
\E_{\tau^\star}\Bigl[\psi^\star(W,X)\bigl\{e^{-\tau^\star}Y(1)-Y(0)\bigr\}\Bigr]
+\E_{\tau^\star}\Bigl[(W-a^\star(X))\bigl\{Y(0)-\mu_0^\star(X)\bigr\}\Bigr]=0 .
\label{eq:expostsplit}
\end{equation}

The second term of \eqref{eq:expostsplit} equals 0: $W-a^\star(X)$ is a function of $(W,X)$, and Assumption \ref{as:spanmeanindep} holds under $\E_{\tau^\star}$, so
\[
\E_{\tau^\star}\bigl[(W-a^\star(X))\{Y(0)-\mu_0^\star(X)\}\bigr]
=\E_{\tau^\star}\bigl[(W-a^\star(X))\{\E_{\tau^\star}[Y(0)\mid W,X]-\mu_0^\star(X)\}\bigr]=0 .
\]
We now change measure back to $\mathbb P$.  Since $\psi^\star(W,X)=0$ unless $W=1$, $e^{\tau^\star W}\psi^\star(W,X)=e^{\tau^\star}\psi^\star(W,X)$, so \eqref{eq:tilt} gives
\[
\E_{\tau^\star}\Bigl[\psi^\star(W,X)\bigl\{e^{-\tau^\star}Y(1)-Y(0)\bigr\}\Bigr]
=
\frac{1}{\E[e^{\tau^\star W}]}\,\E\Bigl[\psi^\star(W,X)\bigl\{Y(1)-e^{\tau^\star}Y(0)\bigr\}\Bigr] .
\]
By \eqref{eq:expostsplit} the left-hand side equals 0, which gives \eqref{eq:expostfinitemoment}.
\end{proof}

\section{Continuous Treatment}

Let $(Y_i,W_i,X_i)_{i=1}^n$ be an i.i.d. sample with $W_i\in[0,\infty)$, potential outcomes $Y_i(w)\ge 0$ defined for all $w\in[0,\infty)$, and $Y_i=Y_i(W_i)$.

\begin{assumption}[Potential outcome paths]
\label{as:contreg}
The data satisfy:
\begin{enumerate}[label=(\roman*)]
    \item \textbf{Tilted moments:} $\E[e^{\tau W}]<\infty$ for every $\tau\in\R$.
    \item \textbf{Absolute continuity:} $w\mapsto Y(w)$ is absolutely continuous on compact subsets of $[0,\infty)$.
    \item \textbf{Integrability:} for every $\tau\in\R$,
    \[
    \E\left[e^{\tau W}\left\{Y(0)+\int_0^\infty e^{-\tau w}\bigl|\partial_wY(w)-\tau Y(w)\bigr|\,dw\right\}^{2}\right]<\infty .
    \]
\end{enumerate}
\end{assumption}

Let $\mathbb P$ be the law of $(Y(\cdot),W,X)$ and let $\mathbb P_\tau$ and $\E_\tau$ be as in \eqref{eq:tilt}, which is well defined by Assumption \ref{as:contreg}(i).  Under $\mathbb P_\tau$ the conditional law of $W$ given $X$ is
\begin{equation}
d\mathbb P_\tau(w\mid X)=\frac{e^{\tau w}\,d\mathbb P(w\mid X)}{C_\tau(X)},
\qquad
C_\tau(X):=\E[e^{\tau W}\mid X],
\label{eq:contcondtilt}
\end{equation}
extending \eqref{eq:condtilt}.  The density in \eqref{eq:tilt} depends on $W$ only, so $\mathbb P_\tau$ leaves the conditional law of $Y(\cdot)$ given $(W,X)$ unchanged; consequently Assumption \ref{as:contdesign} and Assumption \ref{as:spanmeanindep} hold under $\E_\tau$ whenever they hold under $\E$.

As in Section \ref{sec:setup}, the estimand $\tau^\star$ is the coefficient on $W$ in the population Poisson regression with log link \eqref{eq:loglink}.

For $\tau\in\R$ let
\begin{equation}
\pi_\tau(X):=\E_\tau[W\mid X]
\label{eq:contdefpi}
\end{equation}
be the tilted treatment mean, extending \eqref{eq:defpi}.

The purpose of this section is to find conditions, under which we can interpret the same estimand as an average of the potential outcome paths, weighted by the ex-post or the ex-ante weight at level $w\ge 0$.

\begin{definition}[Ex-post weight]
\label{def:contexpost}
\begin{equation}
\psi_\tau(w;W,X):=\bigl(W-\pi_\tau(X)\bigr)\1\{W>w\} .
\label{eq:contexpost}
\end{equation}
\end{definition}

\begin{definition}[Ex-ante weight]
\label{def:contexante}
\begin{equation}
\phi_\tau(w\mid X):=\E_\tau\bigl[\psi_\tau(w;W,X)\mid X\bigr]
=\operatorname{Cov}_\tau\bigl(W,\1\{W>w\}\bigm| X\bigr).
\label{eq:contweights}
\end{equation}
\end{definition}

By such an interpretation we mean that $\tau^\star$ solving the population Poisson first-order conditions \eqref{eq:contcondFOC}--\eqref{eq:conttreatFOC} satisfies
\begin{equation}
\tau^\star
=
\frac{\E_{\tau^\star}\bigl[\int_0^\infty \omega\,e^{-\tau^\star w}\,\partial_wY(w)\,dw\bigr]}
     {\E_{\tau^\star}\bigl[\int_0^\infty \omega\,e^{-\tau^\star w}Y(w)\,dw\bigr]},
\qquad
\omega\in\bigl\{\psi_{\tau^\star}(w;W,X),\ \phi_{\tau^\star}(w\mid X)\bigr\}.
\label{eq:conttarget}
\end{equation}

\paragraph{Reduction to the binary case}

Here we verify that when a binary case $W\in\{0,1\}$ is used, the weights and the estimand reduce to the corresponding definitions and results in Section \ref{sec:setup}.

Specifically, let $W\in\{0,1\}$ and let $Y(\cdot)$ be any absolutely continuous path on $[0,1]$ with endpoints $Y(0)$ and $Y(1)$. We verify that Definitions \ref{def:contexpost} and \ref{def:contexante} reduce to Definitions \ref{def:expost} and \ref{def:exante}, and \eqref{eq:conttarget} to \eqref{eq:expostratio}.

For $0\le w<1$,
\[
\psi_\tau(w;W,X)=\bigl\{1-\pi_\tau(X)\bigr\}W,
\qquad
\phi_\tau(w\mid X)=\pi_\tau(X)\bigl\{1-\pi_\tau(X)\bigr\}=\frac{e^{\tau}\phi_\tau(X)}{C_\tau(X)},
\]
the first being $\psi_\tau(W,X)$ of Definition \ref{def:expost}; for $w\ge 1$ both are 0, since $\1\{W>w\}=0$.

By \eqref{eq:contg}, $\int_0^1 e^{-\tau w}\{\partial_wY(w)-\tau Y(w)\}\,dw=g_\tau(1)-g_\tau(0)=e^{-\tau}Y(1)-Y(0)$, so \eqref{eq:contpseudotrue} is
\[
\E_{\tau^\star}\left[\frac{\phi_{\tau^\star}(X)}{C_{\tau^\star}(X)}\bigl\{Y(1)-e^{\tau^\star}Y(0)\bigr\}\right]=0 .
\]
Under Assumption \ref{as:contdesign}, $\E[e^{\tau W}\mid Y(\cdot),X]=C_\tau(X)$, so by \eqref{eq:tilt} the left-hand side equals $\E\bigl[\phi_{\tau^\star}(X)\{Y(1)-e^{\tau^\star}Y(0)\}\bigr]/\E[e^{\tau^\star W}]$, which is \eqref{eq:pseudotrue}.

\subsection{Useful lemmas}

For $\tau\in\R$ let
\begin{equation}
g_\tau(w):=Y(w)e^{-\tau w},
\qquad
\bar S_\tau:=Y(0)+\int_0^\infty|g_\tau'(w)|\,dw .
\label{eq:contg}
\end{equation}
By Assumption \ref{as:contreg}(ii), $g_\tau'(w)=e^{-\tau w}\bigl\{\partial_wY(w)-\tau Y(w)\bigr\}$ and
\begin{equation}
g_\tau(W)=g_\tau(0)+\int_0^\infty \1\{W>w\}\,g_\tau'(w)\,dw ,
\label{eq:contFTC}
\end{equation}
so $Y=e^{\tau W}g_\tau(W)$, $g_\tau(0)=Y(0)$ and $0\le g_\tau(w)\le\bar S_\tau$ for every $w\ge 0$.
\begin{proposition}
\label{prop:contnonneg}
Suppose Assumption \ref{as:contreg}(i) holds.  For every $\tau\in\R$,
\begin{equation}
\phi_\tau(w\mid X)\ge 0
\qquad\text{for all } w\ge 0 .
\label{eq:contnonneg}
\end{equation}
\end{proposition}

\begin{proof}
By \eqref{eq:contexpost} and \eqref{eq:contweights},
$\phi_\tau(w\mid X)=\E_\tau\bigl[(W-\pi_\tau(X))\1\{W>w\}\bigm| X\bigr]$.

Let $w\ge\pi_\tau(X)$.  On $\{W>w\}$ we have $W>w\ge\pi_\tau(X)$, so $(W-\pi_\tau(X))\1\{W>w\}\ge 0$ and $\phi_\tau(w\mid X)\ge 0$.

Let $w<\pi_\tau(X)$.  Since $\E_\tau[W-\pi_\tau(X)\mid X]=0$ by \eqref{eq:contdefpi}, splitting $1=\1\{W>w\}+\1\{W\le w\}$ gives
\[
\phi_\tau(w\mid X)=-\E_\tau\bigl[(W-\pi_\tau(X))\1\{W\le w\}\bigm| X\bigr].
\]
On $\{W\le w\}$ we have $W\le w<\pi_\tau(X)$, so $(W-\pi_\tau(X))\1\{W\le w\}\le 0$ and $\phi_\tau(w\mid X)\ge 0$.
\end{proof}

\begin{lemma}
\label{lem:contmoments}
Suppose Assumption \ref{as:contreg} holds and let $\tau\in\R$.  Then for every integer $k\ge 0$,
\begin{equation}
\E_\tau[W^k]<\infty,
\qquad
\E_\tau\bigl[\bar S_\tau^{\,2}\bigr]<\infty,
\qquad
\E_\tau\bigl[|W-\pi_\tau(X)|\,\bar S_\tau\bigr]<\infty,
\label{eq:contmoments}
\end{equation}
\begin{equation}
\E\bigl[|W-\pi_\tau(X)|\,Y\bigr]<\infty,
\qquad
\E\bigl[\pi_\tau(X)\,Y\bigr]<\infty,
\label{eq:contYintegrable}
\end{equation}
and
\begin{equation}
\int_0^\infty \phi_\tau(w\mid X)\,dw=\operatorname{Var}_\tau(W\mid X),
\qquad
\E_\tau\left[\int_0^\infty \phi_\tau(w\mid X)\bigl\{e^{-\tau w}Y(w)+|g_\tau'(w)|\bigr\}\,dw\right]<\infty .
\label{eq:contweightmass}
\end{equation}
\end{lemma}

\begin{proof}
For $w\ge 0$, $w^ke^{\tau w}\le c_ke^{(\tau+1)w}$ with $c_k:=\sup_{w\ge 0}w^ke^{-w}<\infty$, so $\E[W^ke^{\tau W}]<\infty$ by Assumption \ref{as:contreg}(i) and $\E_\tau[W^k]<\infty$ by \eqref{eq:tilt}.  Assumption \ref{as:contreg}(iii) and \eqref{eq:tilt} give $\E_\tau[\bar S_\tau^{\,2}]<\infty$.  Since $\pi_\tau(X)=\E_\tau[W\mid X]$,
\[
\E_\tau\bigl[(W-\pi_\tau(X))^2\bigr]
=\E_\tau[W^2]-\E_\tau\bigl[\pi_\tau(X)^2\bigr]
\le\E_\tau[W^2],
\]
so Cauchy--Schwarz gives the third bound in \eqref{eq:contmoments} and $\E_\tau[\pi_\tau(X)\bar S_\tau]<\infty$.  Since $0\le g_\tau(W)\le\bar S_\tau$ and $Y=e^{\tau W}g_\tau(W)$ by \eqref{eq:contg}, applying \eqref{eq:tilt} to these two bounds gives \eqref{eq:contYintegrable}.

Let $A_\tau(X):=\E_\tau[|W-\pi_\tau(X)|\mid X]$, so that $0\le\phi_\tau(w\mid X)\le A_\tau(X)$ for every $w\ge 0$ by \eqref{eq:contweights} and \eqref{eq:contnonneg}.  Since $\E_\tau[W^2\mid X]<\infty$ almost surely, Fubini gives
\[
\int_0^\infty \phi_\tau(w\mid X)\,dw
=\E_\tau\Big[(W-\pi_\tau(X))\int_0^\infty\1\{W>w\}\,dw\;\Big|\;X\Big]
=\operatorname{Var}_\tau(W\mid X).
\]
Hence, using $0\le g_\tau(w)\le\bar S_\tau$,
\[
\int_0^\infty \phi_\tau(w\mid X)\,e^{-\tau w}Y(w)\,dw\le\bar S_\tau\operatorname{Var}_\tau(W\mid X),
\qquad
\int_0^\infty \phi_\tau(w\mid X)\,|g_\tau'(w)|\,dw\le A_\tau(X)\,\bar S_\tau .
\]
By Jensen, $\operatorname{Var}_\tau(W\mid X)^2\le\E_\tau[W^4\mid X]$ and $A_\tau(X)^2\le\E_\tau[(W-\pi_\tau(X))^2\mid X]$, so Cauchy--Schwarz bounds the $\E_\tau$ of both right-hand sides.
\end{proof}

\subsection{Design-Based Assumption}

\begin{assumption}[Design-based]
\label{as:contdesign}
$Y(\cdot)\perp W\mid X$.
\end{assumption}

\begin{proposition}
\label{prop:continuous}
Suppose Assumptions \ref{as:contreg} and \ref{as:contdesign} hold.  Let $\tau^\star\in\R$ and a measurable function $\mu_0^\star(X)\ge 0$ satisfy the population first-order conditions
\begin{align}
\E\big[Y-\mu_0^\star(X)e^{\tau^\star W}\mid X\big] &= 0 \quad \text{a.s.},
\label{eq:contcondFOC}
\\
\E\big[W\{Y-\mu_0^\star(X)e^{\tau^\star W}\}\big] &= 0.
\label{eq:conttreatFOC}
\end{align}
Then
\begin{equation}
\E_{\tau^\star}\left[\int_0^\infty \phi_{\tau^\star}(w\mid X)\,e^{-\tau^\star w}\bigl\{\partial_wY(w)-\tau^\star Y(w)\bigr\}\,dw\right]=0 .
\label{eq:contpseudotrue}
\end{equation}
If in addition $\E_{\tau^\star}\bigl[\int_0^\infty\phi_{\tau^\star}(w\mid X)e^{-\tau^\star w}Y(w)\,dw\bigr]>0$, then
\begin{equation}
\tau^\star
=
\frac{\E_{\tau^\star}\bigl[\int_0^\infty\phi_{\tau^\star}(w\mid X)e^{-\tau^\star w}\,\partial_wY(w)\,dw\bigr]}
     {\E_{\tau^\star}\bigl[\int_0^\infty\phi_{\tau^\star}(w\mid X)e^{-\tau^\star w}Y(w)\,dw\bigr]} ,
\label{eq:contratio}
\end{equation}
\end{proposition}

\begin{proof}
By \eqref{eq:contcondFOC} and \eqref{eq:contcondtilt}, $\E[Y\mid X]=\mu_0^\star(X)C_{\tau^\star}(X)$ and $\E[We^{\tau^\star W}\mid X]=\pi_{\tau^\star}(X)C_{\tau^\star}(X)$, so
\[
\E\bigl[\pi_{\tau^\star}(X)\,\mu_0^\star(X)e^{\tau^\star W}\bigr]
=\E\bigl[W\,\mu_0^\star(X)e^{\tau^\star W}\bigr]
=\E\bigl[\pi_{\tau^\star}(X)\,Y\bigr]<\infty
\]
by \eqref{eq:contYintegrable}.  Together with \eqref{eq:contYintegrable}, every expectation in the next three displays is absolutely convergent.  Multiplying \eqref{eq:contcondFOC} by $\pi_{\tau^\star}(X)$ and using the law of iterated expectations gives
\[
\E\Big[\pi_{\tau^\star}(X)\bigl\{Y-\mu_0^\star(X)e^{\tau^\star W}\bigr\}\Big]=0 .
\]
Subtracting this from \eqref{eq:conttreatFOC},
\[
\E\Big[(W-\pi_{\tau^\star}(X))\bigl\{Y-\mu_0^\star(X)e^{\tau^\star W}\bigr\}\Big]=0 .
\]
By \eqref{eq:tilt} and \eqref{eq:contdefpi},
\[
\E\Big[(W-\pi_{\tau^\star}(X))\,\mu_0^\star(X)e^{\tau^\star W}\Big]
=\E[e^{\tau^\star W}]\,\E_{\tau^\star}\Big[\mu_0^\star(X)\,\E_{\tau^\star}\bigl[W-\pi_{\tau^\star}(X)\bigm| X\bigr]\Big]=0 ,
\]
so
\begin{equation}
\E\bigl[(W-\pi_{\tau^\star}(X))\,Y\bigr]=0 .
\label{eq:contresidualized}
\end{equation}

Since $Y=e^{\tau^\star W}g_{\tau^\star}(W)$, applying \eqref{eq:tilt} to \eqref{eq:contresidualized} and dividing by $\E[e^{\tau^\star W}]>0$,
\[
\E_{\tau^\star}\bigl[(W-\pi_{\tau^\star}(X))\,g_{\tau^\star}(W)\bigr]=0 .
\]
Multiplying \eqref{eq:contFTC} by $W-\pi_{\tau^\star}(X)$ and using $g_{\tau^\star}(0)=Y(0)$ and \eqref{eq:contexpost},
\begin{equation}
(W-\pi_{\tau^\star}(X))\,g_{\tau^\star}(W)
=
Y(0)\,(W-\pi_{\tau^\star}(X))
+
\int_0^\infty \psi_{\tau^\star}(w;W,X)\,g_{\tau^\star}'(w)\,dw .
\label{eq:contidentity}
\end{equation}
Since $|\psi_{\tau^\star}(w;W,X)|\le|W-\pi_{\tau^\star}(X)|$ by \eqref{eq:contexpost}, both terms on the right of \eqref{eq:contidentity} are dominated in absolute value by $|W-\pi_{\tau^\star}(X)|\bar S_{\tau^\star}$, which is integrable by \eqref{eq:contmoments}.

Assumption \ref{as:contdesign} also implies the following conditional independencies:
\[
\E_{\tau^\star}\Big[Y(0)\,(W-\pi_{\tau^\star}(X))\;\Big|\;Y(\cdot),X\Big]
=Y(0)\,\E_{\tau^\star}\bigl[W-\pi_{\tau^\star}(X)\bigm| X\bigr]=0 ,
\]
and, by \eqref{eq:contmoments} and Fubini,
\[
\E_{\tau^\star}\left[\int_0^\infty \psi_{\tau^\star}(w;W,X)\,g_{\tau^\star}'(w)\,dw\;\Bigg|\;Y(\cdot),X\right]
=\int_0^\infty \phi_{\tau^\star}(w\mid X)\,g_{\tau^\star}'(w)\,dw .
\]
Taking $\E_{\tau^\star}$ in \eqref{eq:contidentity}, which is permitted by \eqref{eq:contmoments}, gives \eqref{eq:contpseudotrue}.

Since $e^{-\tau^\star w}\partial_wY(w)=g_{\tau^\star}'(w)+\tau^\star e^{-\tau^\star w}Y(w)$, \eqref{eq:contweightmass} makes the two terms of \eqref{eq:contpseudotrue} separately absolutely convergent, and \eqref{eq:contratio} follows.
\end{proof}

\begin{assumptionprime}[Design-based, restricted controls]
\label{as:contdesignspan}
\begin{equation}
\pi_{\tau^\star}(X)\in\mathrm{span}(R(X))
\qquad\text{almost surely.}
\label{eq:contspancondition}
\end{equation}
\end{assumptionprime}

\begin{corollary}
\label{cor:contfinite}
Suppose Assumptions \ref{as:contreg}, \ref{as:contdesign} and \ref{as:contdesignspan} hold and let $(\beta^\star,\tau^\star)$ solve the population Poisson moment conditions
\begin{equation}
\E\Big[R(X)\bigl\{Y-\mu_0^\star(X)e^{\tau^\star W}\bigr\}\Big]=0,
\qquad
\E\Big[W\bigl\{Y-\mu_0^\star(X)e^{\tau^\star W}\bigr\}\Big]=0,
\label{eq:contfiniteFOC}
\end{equation}
where $\mu_0^\star(X):=\exp\{R(X)'\beta^\star\}$.  Then the conclusions of Proposition \ref{prop:continuous} hold.
\end{corollary}

\begin{proof}
By Assumption \ref{as:contdesignspan} there is $\lambda^\star\in\R^K$ with $\pi_{\tau^\star}(X)=R(X)'\lambda^\star$ almost surely, so the first condition in \eqref{eq:contfiniteFOC} gives
\[
\E\Big[\pi_{\tau^\star}(X)\bigl\{Y-\mu_0^\star(X)e^{\tau^\star W}\bigr\}\Big]
=\lambda^{\star\prime}\,\E\Big[R(X)\bigl\{Y-\mu_0^\star(X)e^{\tau^\star W}\bigr\}\Big]=0 .
\]
Subtracting this from the second condition in \eqref{eq:contfiniteFOC},
\[
\E\Big[(W-\pi_{\tau^\star}(X))\bigl\{Y-\mu_0^\star(X)e^{\tau^\star W}\bigr\}\Big]=0 ,
\]
and the proof of Proposition \ref{prop:continuous} applies from this point unchanged.
\end{proof}

\subsection{Model-Based Assumption}

\begin{assumption}[Tilted linear projection of the treatment]
\label{as:conttiltproj}
$a^\star\in\mathrm{span}(R(X))$ satisfies
\begin{equation}
\E_{\tau^\star}\Big[R(X)\bigl\{W-a^\star(X)\bigr\}\Big]=0 .
\label{eq:conttiltedprojection}
\end{equation}
\end{assumption}

\begin{proposition}
\label{prop:contexpost}
Suppose Assumptions \ref{as:contreg}, \ref{as:spanmeanindep} and \ref{as:conttiltproj} hold, and let $(\beta^\star,\tau^\star)$ solve \eqref{eq:contfiniteFOC} with $\mu_0^\star(X)=\exp\{R(X)'\beta^\star\}$.  Let $\psi^\star(w;W,X):=\bigl(W-a^\star(X)\bigr)\1\{W>w\}$ and suppose
\begin{equation}
\E_{\tau^\star}\bigl[a^\star(X)^2+\mu_0^\star(X)^2\bigr]<\infty .
\label{eq:contexpostintegrable}
\end{equation}
Then
\begin{equation}
\E_{\tau^\star}\left[\int_0^\infty \psi^\star(w;W,X)\,e^{-\tau^\star w}\bigl\{\partial_wY(w)-\tau^\star Y(w)\bigr\}\,dw\right]=0 .
\label{eq:contexpostmoment}
\end{equation}
If in addition $\E_{\tau^\star}\bigl[\int_0^\infty\psi^\star(w;W,X)e^{-\tau^\star w}Y(w)\,dw\bigr]>0$, then
\begin{equation}
\tau^\star
=
\frac{\E_{\tau^\star}\bigl[\int_0^\infty\psi^\star(w;W,X)e^{-\tau^\star w}\,\partial_wY(w)\,dw\bigr]}
     {\E_{\tau^\star}\bigl[\int_0^\infty\psi^\star(w;W,X)e^{-\tau^\star w}Y(w)\,dw\bigr]} .
\label{eq:contexpostratio}
\end{equation}
\end{proposition}

\begin{proof}
Since $a^\star=R(X)'\lambda^\star$ for some $\lambda^\star\in\R^K$, the first condition in \eqref{eq:contfiniteFOC} gives
\[
\E\Big[a^\star(X)\bigl\{Y-\mu_0^\star(X)e^{\tau^\star W}\bigr\}\Big]
=\lambda^{\star\prime}\,\E\Big[R(X)\bigl\{Y-\mu_0^\star(X)e^{\tau^\star W}\bigr\}\Big]=0 .
\]
Subtracting this from the second condition in \eqref{eq:contfiniteFOC},
\begin{equation}
\E\Big[(W-a^\star(X))\bigl\{Y-\mu_0^\star(X)e^{\tau^\star W}\bigr\}\Big]=0 .
\label{eq:contexpostsubtraction}
\end{equation}

By \eqref{eq:contg},
\[
Y-\mu_0^\star(X)e^{\tau^\star W}
=
e^{\tau^\star W}\bigl\{g_{\tau^\star}(W)-g_{\tau^\star}(0)\bigr\}+e^{\tau^\star W}\bigl\{Y(0)-\mu_0^\star(X)\bigr\},
\]
so multiplying by $W-a^\star(X)$ and applying \eqref{eq:tilt} to \eqref{eq:contexpostsubtraction}, then dividing by $\E[e^{\tau^\star W}]>0$,
\begin{equation}
\E_{\tau^\star}\Bigl[(W-a^\star(X))\bigl\{g_{\tau^\star}(W)-g_{\tau^\star}(0)\bigr\}\Bigr]
+\E_{\tau^\star}\Bigl[(W-a^\star(X))\bigl\{Y(0)-\mu_0^\star(X)\bigr\}\Bigr]=0 .
\label{eq:contexpostidentity}
\end{equation}
The split is permitted because both terms are absolutely integrable: by \eqref{eq:contexpostintegrable} and \eqref{eq:contmoments},
\[
\E_{\tau^\star}\bigl[(W-a^\star(X))^2\bigr]<\infty,
\qquad
\E_{\tau^\star}\Bigl[\bigl\{\bar S_{\tau^\star}+\mu_0^\star(X)\bigr\}^2\Bigr]<\infty ,
\]
and Cauchy--Schwarz bounds the product, using $Y(0)\le\bar S_{\tau^\star}$ and $|g_{\tau^\star}(W)-g_{\tau^\star}(0)|\le\bar S_{\tau^\star}$ from \eqref{eq:contg}.

The second term of \eqref{eq:contexpostidentity} equals 0: $W-a^\star(X)$ is a function of $(W,X)$, and Assumption \ref{as:spanmeanindep} holds under $\E_{\tau^\star}$, so
\[
\E_{\tau^\star}\bigl[(W-a^\star(X))\{Y(0)-\mu_0^\star(X)\}\bigr]
=\E_{\tau^\star}\bigl[(W-a^\star(X))\{\E_{\tau^\star}[Y(0)\mid W,X]-\mu_0^\star(X)\}\bigr]=0 .
\]
In the first term, \eqref{eq:contFTC} gives $(W-a^\star(X))\{g_{\tau^\star}(W)-g_{\tau^\star}(0)\}=\int_0^\infty \psi^\star(w;W,X)g_{\tau^\star}'(w)\,dw$, and the same bound permits interchanging $\E_{\tau^\star}$ with $\int_0^\infty dw$.  This gives \eqref{eq:contexpostmoment}, and \eqref{eq:contexpostratio} is \eqref{eq:contexpostmoment} rearranged.
\end{proof}

\begin{proposition}
\label{prop:contmonotone}
Suppose Assumptions \ref{as:contdesign} and \ref{as:contreg}(i) hold, $\E[Y]<\infty$, and set
\[
M(\tau):=\E\bigl[(W-\pi_\tau(X))\,Y\bigr],
\qquad
V_\tau(X):=\operatorname{Var}_\tau\bigl(W\bigm| X\bigr).
\]
Fix $\tau\in\R$ and suppose $\E\bigl[\sup_{|t-\tau|\le\delta}V_t(X)\,Y\bigr]<\infty$ for some $\delta>0$.  Then
\[
M'(\tau)=-\E\bigl[V_\tau(X)\,\E[Y\mid X]\bigr]\le 0 ,
\]
with equality if and only if $V_\tau(X)\,\E[Y\mid X]=0$ almost surely.
\end{proposition}

\begin{proof}
By \eqref{eq:contcondtilt}, $\pi_\tau(X)=\E[We^{\tau W}\mid X]/C_\tau(X)$ and $\E_\tau[W^2\mid X]=\E[W^2e^{\tau W}\mid X]/C_\tau(X)$.  Differentiating the first in $\tau$, which is justified by Assumption \ref{as:contreg}(i),
\[
\partial_\tau\pi_\tau(X)
=\frac{\E[W^2e^{\tau W}\mid X]}{C_\tau(X)}-\pi_\tau(X)^2
=\E_\tau[W^2\mid X]-\E_\tau[W\mid X]^2
=V_\tau(X).
\]
Since $Y$ does not depend on $\tau$, the domination hypothesis gives
\[
M'(\tau)=-\E\bigl[\partial_\tau\pi_\tau(X)\,Y\bigr]=-\E\bigl[V_\tau(X)\,Y\bigr]
=-\E\bigl[V_\tau(X)\,\E[Y\mid X]\bigr],
\]
which is nonpositive because $V_\tau(X)\ge 0$ and $Y\ge 0$.
\end{proof}

\section{Useful Facts}

For $W\in\{0,1\}$, with $\pi_\tau(X)$ and $\phi_\tau(X)$ as in \eqref{eq:defpi} and \eqref{eq:defexante}, on $\{0<e(X)<1\}$,
\begin{equation}
\operatorname{logit}\pi_\tau(X)=\operatorname{logit}e(X)+\tau,
\qquad
\frac{1}{\phi_\tau(X)}=\frac{1}{e(X)}+\frac{e^{\tau}}{1-e(X)}.
\label{eq:logitform}
\end{equation}
For $W\in[0,\infty)$, with $\pi_\tau(X)$ and $C_\tau(X)$ as in \eqref{eq:contdefpi} and \eqref{eq:contcondtilt},
\[
\pi_\tau(X)=\partial_\tau\log C_\tau(X).
\]

\end{document}

%%% Local Variables:
%%% mode: latex
%%% TeX-master: t
%%% End:
